\documentclass{article}
\usepackage{fleqn}
\usepackage{epsf}
\usepackage{aima2e-slides}

\begin{document}

\begin{huge}
\titleslide{Các quyết định phức tạp (Complex decisions)}{Chương 17, Phần 1--3}

\sf

%%%%%%%%%%%% Slide %%%%%%%%%%%%%%%%%%%%%%%%%%%%%%%%%%%%%%%%%%%%%%%%%%%
\heading{Phác thảo}

\blob Các vấn đề về quyết định tuần tự

\blob Lặp lại giá trị

\blob Lặp lại chính sách


%%%%%%%%%%%% Slide %%%%%%%%%%%%%%%%%%%%%%%%%%%%%%%%%%%%%%%%%%%%%%%%%%%
\heading{Các vấn đề về quyết định tuần tự}

\vspace*{0.3in}

\epsfxsize=1.06\textwidth
\fig{\file{figures}{decision-problems.ps}}



%%%%%%%%%%%% Slide %%%%%%%%%%%%%%%%%%%%%%%%%%%%%%%%%%%%%%%%%%%%%%%%%%%
\heading{Ví dụ MDP}

\epsffile{\file{figures}{sequential-decision-world.ps}}

\centerline{%
\epsfxsize=0,4\textwidth
\epsffile{\file{figures}{sequential-decision-world.ps}} \hspace*{1in} $s \in S$
\epsfxsize=0,2\textwidth
$a \in A$}

Kỳ $s \in S$, hành động $a \in A$

\u{Mô hình} $T(s,a,s') \equiv P(s'|s,a)$ = xác suất $a$ trong $s$ dẫn đến $s'$

\u{Chức năng thưởng} $R(s)$ (hoặc $R(s,a)$, $R(s,a,s')$)\nl
   = $\left\{\begin{array}{ll} -0.04 & \mbox{ (small penalty) for nonterminal states}\\
                               \pm 1 & \mbox{ for terminal states}\end{array}\right.$


%%%%%%%%%%%% Slide %%%%%%%%%%%%%%%%%%%%%%%%%%%%%%%%%%%%%%%%%%%%%%%%%%%
\heading{Giải quyết MDP}

Trong các bài toán tìm kiếm, mục đích là tìm một chuỗi {\em} tối ưu

Trong MDP, mục tiêu là tìm ra một {\em Policy} tối ưu $\pi(s)$\nl
   tức là hành động tốt nhất cho mọi trạng thái có thể $s$ \nl
   (vì không thể đoán trước được mình sẽ đi về đâu)\\
Chính sách tối ưu tối đa hóa (giả sử) {\em tổng phần thưởng dự kiến}

Chính sách tối ưu khi hình phạt của tiểu bang $R(s)$ là --0,04:

\vspace*{0.2in}

\epsfxsize=0,45\textwidth
\fig{\file{figures}{sequential-decision-policy.ps}}

%%%%%%%%%%%% Slide %%%%%%%%%%%%%%%%%%%%%%%%%%%%%%%%%%%%%%%%%%%%%%%%%%%
\heading{Rủi ro và phần thưởng}

\vspace*{0.2in}

\epsfxsize=0,75\textwidth
\fig{\file{figures}{policy-flips4.ps}}



%%%%%%%%%%%% Slide %%%%%%%%%%%%%%%%%%%%%%%%%%%%%%%%%%%%%%%%%%%%%%%%%%%
\heading{Tiện ích của chuỗi trạng thái}

Cần hiểu các ưu tiên giữa {\em chuỗi} trạng thái

Thường xem xét \u{tùy chọn cố định} trên chuỗi phần thưởng:
\[
  [r,r_0,r_1,r_2,\ldots]\pref [r,r'_0,r'_1,r'_2,\ldots]
  \lequiv
  [r_0,r_1,r_2,\ldots]\pref [r'_0,r'_1,r'_2,\ldots]
\]  
\u{Định lý}: chỉ có hai cách để kết hợp phần thưởng theo thời gian.\al
  1) Hàm tiện ích {\em Additive}:\nl
   $U([s_0,s_1,s_2,\ldots]) = R(s_0) + R(s_1) + R(s_2) + \cdots $\al
  2) Hàm tiện ích {\em Giảm giá}:\nl
   $U([s_0,s_1,s_2,\ldots]) = R(s_0) + \gamma R(s_1) + \gamma^2 R(s_2) + \cdots $\nl
   trong đó $\gamma$ là \u{hệ số chiết khấu}

%%%%%%%%%%%% Slide %%%%%%%%%%%%%%%%%%%%%%%%%%%%%%%%%%%%%%%%%%%%%%%%%%%
\heading{Tiện ích của các trạng thái}

Tiện ích của một {\em state} (hay còn gọi là {\em value}) được xác định là \nl
$U(s) = {}$ \u{tổng số phần thưởng dự kiến (được chiết khấu) (cho đến khi chấm dứt)}\nl
\phantom{$U(s) = {}$} \u{giả sử hành động tối ưu}

Với những tiện ích của các bang, việc lựa chọn hành động tốt nhất chỉ là MEU:\al
  tối đa hóa lợi ích mong đợi của những người kế nhiệm trực tiếp 

\vspace*{0.2in}

\twofig{\file{figures}{sequential-decision-values.ps}}{\file{figures}{sequential-decision-policy.ps}}


%%%%%%%%%%%% Slide %%%%%%%%%%%%%%%%%%%%%%%%%%%%%%%%%%%%%%%%%%%%%%%%%%%
\heading{Tiện ích tiếp theo.}

Vấn đề: các tiện ích cộng thêm có thời gian sống vô hạn $\implies$ là vô hạn

1) \u{Chân trời hữu hạn}: kết thúc tại một {\em thời gian cố định} $T$\al
  Chính sách $\implies$ \u{không cố định}: $\pi(s)$ phụ thuộc vào thời gian còn lại

2) \u{(Các) trạng thái hấp thụ}: w/ prob.~1, tác nhân cuối cùng `` chết '' đối với bất kỳ $\pi$\al
  $\implies$ Tiện ích dự kiến của mọi trạng thái là hữu hạn

3) Chiết khấu: giả sử $\gamma<1$, $R(s) \leq R_{{\rm max}}$,
\[ U([s_0,\ldots s_{\infty}]) = \mysum_{t=0}^{\infty} \gamma^t R(s_t) \leq R_{{\rm max}}/(1-\gamma)\]
Đường chân trời ngắn hơn $\gamma \Rightarrow {}$ nhỏ hơn

4) Tối đa hóa \u{lợi ích hệ thống} = phần thưởng trung bình mỗi bước thời gian\\
Định lý: chính sách tối ưu có mức tăng không đổi sau thoáng qua ban đầu\\
Ví dụ: kế hoạch hàng ngày của tài xế taxi đưa đón hành khách



%%%%%%%%%%%% Slide %%%%%%%%%%%%%%%%%%%%%%%%%%%%%%%%%%%%%%%%%%%%%%%%%%%
\heading{Lập trình động: phương trình Bellman}

Định nghĩa về tính hữu dụng của các trạng thái dẫn đến một mối quan hệ đơn giản giữa
Tiện ích của các nước lân cận:

\u{tổng số phần thưởng dự kiến}\al
= \u{phần thưởng hiện tại}\nl
+ $\gamma \stimes{}$\u{tổng số phần thưởng dự kiến sau khi thực hiện hành động tốt nhất}

Phương trình Bellman (1957):
\[ U(s) = R(s) + \gamma\,\max_a \mysum_{s'} U(s') T(s,a,s')\]

$U(1,1) = -0.04$\\
\tab + $\gamma\,\max\{ 0.8 U(1,2) + 0.1 U(2,1) + 0.1 U(1,1),$\hfill{\em up\\
\tab \phantom{+ \mbox{$\max\{$}}$0.9 U(1,1) + 0.1 U(1,2) $\hfill{\em left\\
\tab \phantom{+ \mbox{$\max\{$}}$0.9 U(1,1) + 0.1 U(2,1) $\hfill{\em down\\
\tab \phantom{+ \mbox{$\max\{$}$0.8 U(2,1) + 0.1 U(1,2) + 0.1 U(1,1) \}$\hfill{\em đúng}

Một phương trình trên mỗi trạng thái = $n$ \u{phi tuyến } phương trình trong $n$ ẩn số


%%%%%%%%%%%% Slide %%%%%%%%%%%%%%%%%%%%%%%%%%%%%%%%%%%%%%%%%%%%%%%%%%%
\heading{Thuật toán lặp giá trị}

\u{Idea}: Bắt đầu với các giá trị tiện ích tùy ý\nl
          Cập nhật để làm cho chúng \u{nhất quán cục bộ} với Bellman eqn.\nl
          Tối ưu toàn cầu nhất quán cục bộ $\Rightarrow$ ở mọi nơi

Lặp lại đồng thời cho mọi $s$ cho đến khi ``không thay đổi''
\[ U(s) \leftarrow R(s) + \gamma\,\max_a \mysum_{s'} U(s') T(s,a,s') \qquad \mbox{for all } s\]

\epsfxsize=0,6\textwidth
\fig{\file{graphs}{4x3-vi-curve.ps}}

%%%%%%%%%%%% Slide %%%%%%%%%%%%%%%%%%%%%%%%%%%%%%%%%%%%%%%%%%%%%%%%%%%
\heading{Hội tụ}

Xác định \u{max-norm} $||U|| = \max_s |U(s)|$,\\
vì vậy $||U-V||$ = chênh lệch tối đa giữa $U$ và $V$

Đặt $U^t$ và $U^{t+1}$ là các giá trị gần đúng liên tiếp của hữu dụng thực $U$

\u{Định lý}: Đối với hai phép tính gần đúng bất kỳ $U^t$ và $V^t$
\[ ||U^{t+1}-V^{t+1}|| \leq \gamma\, ||U^{t}-V^{t}|| \]
Tức là, mọi phép tính gần đúng khác biệt đều phải tiến gần nhau hơn\\
vì vậy, đặc biệt, bất kỳ phép tính gần đúng nào cũng phải tiến gần hơn đến giá trị thực $U$\\
và sự lặp lại giá trị hội tụ thành một giải pháp duy nhất, ổn định, tối ưu

\u{Định lý}: nếu $||U^{t+1} - U^t||<\epsilon$ thì $||U^{t+1} - U||<2\epsilon\gamma/(1-\gamma)$\\
Tức là, khi thay đổi trong $U^t$ trở nên nhỏ, chúng ta gần như đã hoàn tất.

Chính sách MEU sử dụng $U^t$ có thể tối ưu từ lâu trước khi hội tụ các giá trị

%%%%%%%%%%%% Slide %%%%%%%%%%%%%%%%%%%%%%%%%%%%%%%%%%%%%%%%%%%%%%%%%%%
\heading{Lặp lại chính sách}

Howard, 1960: tìm kiếm đồng thời các giá trị chính sách và tiện ích tối ưu

Thuật toán:\al
  $\pi \leftarrow {}$ chính sách ban đầu tùy ý\al
  lặp lại cho đến khi không có thay đổi nào trong $\pi$\nl
    các tiện ích tính toán được cung cấp $\pi$ \nl
    cập nhật $\pi$ như thể các tiện ích đã chính xác (tức là MEU cục bộ)

Để tính toán các tiện ích được xác định giá trị $\pi$ (\u{} cố định):
\[ U(s) = R(s) + \gamma\,\mysum_{s'} U(s') T(s,\pi(s),s')\qquad \mbox{for all } s \]
tức là, $n$ phương trình \u{tuyến tính} đồng thời trong $n$ ẩn số,
giải quyết trong $O(n^3)$

%%%%%%%%%%%% Slide %%%%%%%%%%%%%%%%%%%%%%%%%%%%%%%%%%%%%%%%%%%%%%%%%%%
\heading{ Lặp lại chính sách đã sửa đổi }

Việc lặp lại chính sách thường hội tụ trong một vài lần lặp, nhưng mỗi lần đều tốn kém

Ý tưởng: sử dụng một vài bước lặp giá trị (nhưng đã sửa $\pi$) \\
bắt đầu từ hàm giá trị được tạo lần trước \\
để tạo ra bước xác định giá trị gần đúng.

Thường hội tụ nhanh hơn nhiều so với VI hoặc PI thuần túy

Dẫn đến các thuật toán tổng quát hơn trong đó cập nhật giá trị Bellman
và cập nhật chính sách của Howard có thể được thực hiện cục bộ theo bất kỳ thứ tự nào

Thuật toán \u{Học tăng cường} hoạt động bằng cách thực hiện như vậy
cập nhật dựa trên các chuyển đổi được quan sát được thực hiện theo cách chưa biết ban đầu
môi trường

%%%%%%%%%%%% Slide %%%%%%%%%%%%%%%%%%%%%%%%%%%%%%%%%%%%%%%%%%%%%%%%%%%
\heading{Khả năng quan sát một phần}

POMDP có mô hình quan sát \u{} $O(s,e)$ xác định xác suất mà
đại lý thu được bằng chứng $e$ khi ở trạng thái $s$

Tác nhân không biết nó đang ở trạng thái nào \nl
$\implies$ nói về chính sách thật vô nghĩa $\pi(s)$!!

\u{Định lý} (Astrom, 1965): chính sách tối ưu trong POMDP là một hàm\nl
  $\pi(b)$ trong đó $b$ là \u{trạng thái niềm tin} (phân bố xác suất trên các trạng thái)

Có thể chuyển đổi POMDP thành MDP trong không gian trạng thái niềm tin, trong đó\nl
  $T(b,a,b')$ là xác suất trạng thái niềm tin mới là $b'$\nl
  cho rằng trạng thái niềm tin hiện tại là $b$ và tác nhân thực hiện $a$.\nl
  Tức là về cơ bản là bước cập nhật lọc

%%%%%%%%%%%% Slide %%%%%%%%%%%%%%%%%%%%%%%%%%%%%%%%%%%%%%%%%%%%%%%%%%%
\heading{Tiếp theo khả năng quan sát một phần }

Giải pháp tự động bao gồm hành vi thu thập thông tin

Nếu có trạng thái $n$ thì $b$ là một vectơ có giá trị thực theo chiều $n$\nl
$\implies$ việc giải quyết POMDP rất (thực sự là PSPACE-) rất khó!

Thế giới thực là một POMDP (ban đầu không xác định được $T$ và $O$)




\end{huge}
\end{document}







%%%%%%%%%%%% Slide %%%%%%%%%%%%%%%%%%%%%%%%%%%%%%%%%%%%%%%%%%%%%%%%%%%
\heading{Biến đổi bất biến}

Hành vi không tuần tự bất biến dưới biến đổi tuyến tính:
\[
  U'(s) = k_1 U(s) + k_2 \quad\mbox{where } k_1 > 0
\]
Hành vi tuần tự với tiện ích phụ gia là bất biến\\
bổ sung bất kỳ phần thưởng \u{dựa trên tiềm năng} nào:
\[
  R'(s,a,s') = R(s,a,s') + F(s,a,s')
\]
trong đó phần thưởng được thêm vào phải thỏa mãn
\[
\gamma \Phi(s') - \Phi(s)
\]
đối với một số \u{hàm tiềm năng } $\Phi$ ở các trạng thái


%%%%%%%%%%%% Slide %%%%%%%%%%%%%%%%%%%%%%%%%%%%%%%%%%%%%%%%%%%%%%%%%%%
\heading{Q-lặp }

Xác định $Q(s,a)$ = giá trị mong đợi của việc thực hiện hành động $a$ ở trạng thái $s$\al
  = $R(s) + \gamma\,\mysum_{s'} U(s') T(s,a,s')$\al
  tức là $U(s) = \max_a Q(s,a)$

Thuật toán Q-iteration: giống như phép lặp giá trị, nhưng thực hiện
\[
Q(s,a) \leftarrow R(s) + \gamma\,\mysum_{s'} T(s,a,s') \max_{a'} Q(s',a') \qquad \mbox{for all } s,a
\]


%%%%%%%%%%%% Slide %%%%%%%%%%%%%%%%%%%%%%%%%%%%%%%%%%%%%%%%%%%%%%%%%%%
\heading{Vấn đề}

Độ phức tạp: đa thời gian về số lượng trạng thái (bằng lập trình tuyến tính)\al
  nhưng số lượng trạng thái là số mũ theo số lượng biến trạng thái\al
  $\rightarrow$ Boutilier {\em et al}>: sử dụng cấu trúc trạng thái\al
  $\rightarrow$ học tăng cường: mẫu $S$, xấp xỉ $U$\al
  $\rightarrow$ phương pháp phân cấp trong xây dựng chính sách

Khả năng quan sát một phần: không thể đảm nhận việc quan sát hoàn hảo trạng thái\al
  $\implies$ duy trì \u{trạng thái niềm tin} (trạng thái sau)\nl
  và giảm POMDP xuống MDP trong không gian trạng thái niềm tin (eek)\al
  $\rightarrow$ Cassandra {\em et al}: Lặp lại giá trị POMDP\al
  $\rightarrow$ Hansen: Chính sách POMDP dưới dạng DFA\al

